
\section{What to Evolve?}
\label{sec:what}

A self-evolving agent differs from a static agent not by \emph{what} components it contains,
but by \emph{which internal states} can be autonomously modified based on its own
trajectories, reflections, and feedback signals.  
Thus, the key question of this section is to identify the \textit{evolutionary loci} within an
agent system $\Pi = (\Gamma, \{\psi_i\}, \{C_i\}, \{\mathcal{W}_i\})$—the parts of the system whose
states can be rewritten in an experience-driven and persistent manner, enabling cumulative self-improvement.


Following the formulation in Section~\ref{sec:definition}, 
these evolutionary loci align with four major pillars of an agent system.
Our investigation starts at the agent's cognitive core, namely the \textbf{Models} 
$\{\psi_i\}$,
whose parameters can be continuously updated through self-generated supervision, execution traces, or environmental feedback
\citep{zhou2025selfchallenginglanguagemodelagents, wang2025ragen}. 
We then consider the \textbf{Context} $\{C_i\}$ 
--including instructions
\citep{xiang2025self, khattab2023dspy} 
and long-term memory
\citep{chhikara2025mem0, wang2024agent}
--which evolves as agents reflect, store, and retrieve experience in ways that
shape future decision-making.
From this internal foundation, we 
examine the evolution of
\textbf{Tools} $\{\mathcal{W}_i\}$,
where agents autonomously create
\citep{qiu2025alita}, 
refine
\citep{fromexplorationtomastery}, 
and managing
\citep{wang2025toolgen} 
executable skills based on verifiable interaction signals
Finally, we scale 
to the \textbf{Agentic Architecture}, where the system’s
\textbf{architecture}
\citep{hu2024automated, zhang2024aflow} 
and collaborative structures
\citep{wan2025remalearningmetathinkllms} 
are optimized over time, enabling structural adaptation beyond individual components.
We present 
representative examples of these evolutionary loci in Table~\ref{tab:what_to_evolve}.


\subsection{Models}

Models constitute a primary \emph{locus of self-evolution}, as their parameters can be
autonomously rewritten based on the agent’s own trajectories, reflections, and
interaction outcomes.
The ability of these models to evolve by continually adapting their internal parameters and expanding their functional capabilities
is essential for the development of autonomous, general-purpose agents. 
Unlike static systems that rely heavily on human-annotated datasets and fixed training regimes, 
self-evolving models can improve through interaction, self-supervised data generation, and dynamic learning loops, thereby achieving greater efficiency, adaptability, and scalability.

In detail, we outline the principal axes along which model evolution unfolds. 
These include learning from self-generated supervision to refine model weights, 
evolving through interaction with constructed or external environments, and
integrating feedback signals that directly reshape future reasoning behaviors. Together,
these strategies represent a shift from passive learning paradigms toward active self-improvement.



\begin{table}[!t]
\small
\centering
\caption{Representative self-evolving agent methods positioned along four evolutionary pillars; a filled bullet ($\bullet$) marks dimensions where the approach actively evolves. }
\label{tab:what_to_evolve}
\resizebox{0.9\textwidth}{!}{
\begin{tabular}{@{}l cc cc ccc cc@{}}
\toprule
\textbf{Method} &
\multicolumn{2}{c}{\textbf{Model}} &
\multicolumn{2}{c}{\textbf{Context}} &
\multicolumn{3}{c}{\textbf{Tool}} &
\multicolumn{2}{c@{}}{\textbf{Architecture}} \\
\cmidrule(lr){2-3} \cmidrule(lr){4-5} \cmidrule(lr){6-8} \cmidrule(l){9-10}
 & Policy & Experience & Prompt & Memory & Creation & Mastery & Selection & Single & Multi \\
\midrule
SCA\citep{zhou2025selfchallenginglanguagemodelagents} & \(\bullet\) & \(\bullet\) & \(\circ\) & \(\circ\) &
\(\bullet\) & \(\circ\) & \(\circ\) & \(\circ\) & \(\circ\) \\
RAGEN\citep{wang2025ragen} & \(\bullet\) & \(\bullet\) & \(\bullet\) & \(\circ\) &
\(\circ\) & \(\circ\) & \(\circ\) & \(\bullet\)  & \(\circ\) \\
AgentGen \citep{hu2024agentgen} & \(\circ\) & \(\bullet\) & \(\bullet\) & \(\bullet\)  &
\(\bullet\) & \(\circ\) & \(\circ\) & \(\bullet\)  & \(\circ\) \\
Promptbreeder\citep{fernando2023promptbreeder}& \(\circ\) & \(\circ\) & \(\bullet\)  & \(\circ\) &
\(\circ\) & \(\circ\) & \(\circ\) & \(\bullet\)  & \(\circ\) \\
Expel\citep{zhao2024expel} & \(\circ\) & \(\bullet\)  & \(\circ\) & \(\bullet\) &
\(\circ\) & \(\circ\) & \(\circ\) & \(\circ\) & \(\circ\) \\
Agent Workflow Memory\citep{wang2024agent} & \(\circ\) & \(\circ\) & \(\circ\) &  \(\bullet\) &
\(\circ\) &  \(\circ\) &  \(\bullet\) & \(\circ\) & \(\circ\) \\
Mem0\citep{chhikara2025mem0} & \(\circ\) & \(\circ\) & \(\circ\) & \(\bullet\) &
\(\circ\) & \(\circ\) & \(\circ\) & \(\circ\) & \(\circ\) \\
MAS-Zero\citep{ke2025maszero} & \(\circ\) & \(\circ\) &  \(\bullet\) & \(\circ\) &
\(\circ\) &  \(\circ\) & \(\circ\) & \(\circ\) &  \(\bullet\) \\
Multi-Agent Design\citep{zhou2025maasprompt} & \(\circ\) & \(\circ\) &  \(\bullet\) & \(\circ\) &
\(\circ\) & \(\circ\) &  \(\bullet\) & \(\circ\) &  \(\bullet\) \\
SPO\citep{xiang2025self} & \(\circ\) & \(\circ\) & \(\bullet\) & \(\circ\) &
\(\circ\) & \(\circ\) & \(\circ\) & \(\circ\) & \(\circ\) \\
Alita\citep{qiu2025alita} & \(\circ\) & \(\circ\) & \(\circ\) & \(\circ\) &
\(\bullet\) & \(\circ\) & \(\bullet\) & \(\circ\) & \(\circ\) \\
TextGrad\citep{yellamraju2024textgrad} & \(\circ\) & \(\circ\) & \(\bullet\) & \(\circ\) &
\(\circ\) & \(\bullet\) & \(\bullet\) & \(\bullet\) & \(\circ\) \\
DGM\citep{zhang2025darwin} & \(\circ\) & \(\circ\) & \(\bullet\) & \(\circ\) &
\(\circ\) & \(\circ\) & \(\circ\) & \(\bullet\) & \(\circ\) \\
AlphaEvolve\citep{novikov2025alphaevolve} & \(\circ\) & \(\circ\) & \(\bullet\)  & \(\circ\) &
\(\bullet\)  & \(\bullet\)  & \(\circ\) & \(\bullet\) & \(\circ\) \\
ADAS\citep{hu2024automated} & \(\circ\) & \(\circ\) & \(\bullet\)& \(\circ\) &
\(\bullet\) & \(\circ\) & \(\circ\) & \(\bullet\) & \(\bullet\) \\
AFlow\citep{zhang2024aflow} & \(\circ\) & \(\circ\) & \(\bullet\)& \(\circ\) &
\(\bullet\) & \(\circ\) & \(\bullet\) & \(\bullet\) & \(\bullet\) \\
ReMA\citep{wan2025remalearningmetathinkllms} & \(\circ\) & \(\circ\) & \(\circ\) & \(\circ\) &
\(\circ\) & \(\circ\) & \(\circ\) & \(\circ\) & \(\bullet\) \\
SkillWeaver\citep{zheng2025skillweaver} & \(\circ\) & \(\circ\) & \(\circ\) & \(\bullet\) &
\(\bullet\) & \(\bullet\) & \(\bullet\) & \(\circ\) & \(\circ\) \\
LearnAct\citep{zhao2024learnact} & \(\circ\) & \(\circ\) & \(\circ\) & \(\bullet\) &
\(\bullet\) & \(\circ\) & \(\bullet\) & \(\circ\) & \(\circ\) \\
DRAFT\citep{fromexplorationtomastery} & \(\circ\) & \(\circ\) & \(\bullet\) & \(\circ\) &
\(\bullet\) & \(\bullet\) & \(\circ\) & \(\circ\) & \(\circ\) \\
ToolGen\citep{wang2025toolgen} & \(\circ\) & \(\circ\) & \(\circ\) & \(\bullet\) &
\(\bullet\)& \(\bullet\) & \(\circ\) & \(\circ\) & \(\circ\) \\
CRAFT\citep{yuan2024craft} & \(\circ\) & \(\circ\) & \(\circ\) & \(\circ\) &
\(\bullet\) & \(\bullet\) & \(\bullet\) & \(\circ\) & \(\circ\) \\
CREATOR\citep{qian2023creator} & \(\circ\) & \(\circ\) & \(\circ\) & \(\circ\) &
\(\bullet\) & \(\bullet\) & \(\circ\) & \(\circ\) & \(\circ\) \\
Voyager\citep{wang2023voyager} & \(\circ\) & \(\circ\) & \(\bullet\) & \(\bullet\) &
\(\bullet\) & \(\bullet\) & \(\bullet\) & \(\circ\) & \(\bullet\) \\
\bottomrule
\end{tabular}}
\end{table}



\paragraph{Policy}
A self-evolving agent can refine its parameters to perform better on targeted tasks. Traditional methods of data collection for training agents on tool-use benchmarks are costly and often yield limited coverage, while purely synthetic data-generation pipelines typically suffer from inadequate quality. Consequently, recent studies emphasize enabling agents to autonomously generate data to improve their own model weights.
One representative approach is the Self-Challenging Agent (SCA)\citep{zhou2025selfchallenginglanguagemodelagents}, where a language model alternates roles between a challenger generating executable Code-as-Task problems and an executor solving them. The model then fine-tunes its parameters using trajectories derived from successful solutions, resulting in significant performance gains on complex, multi-step tasks. Similarly, the Self-Rewarding Self-Improving framework\citep{simonds2025selfrewardingselfimproving} implements an internal self-judging mechanism, allowing the model to autonomously generate problems, solve them, and assess its performance, thus producing self-contained fine-tuning data without external annotations. This method demonstrated notable improvements, particularly in complex reasoning tasks.
Beyond task creation, another promising research direction involves leveraging interaction feedback directly for parameter updates. For instance, SELF\citep{lu2023self}, SCoRe\citep{kumar2024training}, and PAG\citep{jiang2025pag} interpret execution traces or natural-language critiques as reward signals within an online Supervised Fine-Tuning (SFT) combined with Reinforcement Learning (RL) framework, enabling continuous policy improvement. TextGrad\citep{yellamraju2024textgrad} further extends this concept by treating unstructured textual feedback as a differentiable training signal capable of directly influencing both prompt design and model parameters. Additionally, AutoRule\citep{wang2025autorule} converts language-model reasoning traces and preference feedback into explicit rule-based training rewards, enhancing the quality of model outputs through structured reward signals. 
Collectively, these advancements chart a clear trajectory—from agents autonomously crafting their training tasks to directly refining their parameters based on execution feedback, highlighting the capacity of models to evolve continuously by learning from the data they produce.
\paragraph{Experience}
Agents can evolve not only by adjusting their internal parameters but also by actively interacting with or even constructing their environments, capturing experiences, and transforming them into learning signals that drive iterative improvement. This environmental loop provides agents with the complexity and diversity required for scalable self-adaptation.
The Self-Challenging Agent (SCA)\citep{zhou2025selfchallenginglanguagemodelagents} exemplifies this dynamic at the task level, where the agent autonomously generates novel Code-as-Task problems, executes them, and then filters successful trajectories for retraining itself. AgentGen\citep{hu2024agentgen} extends this concept to full-environment generation, synthesizing diverse simulation worlds (in PDDL or Gym-style formats) derived from an initial corpus. It implements a bidirectional evolution loop that progressively adjusts task difficulty, enabling the agent to continuously grow within a dynamically structured curriculum. Reflexion\citep{shinn2023reflexion} complements this by introducing self-reflective mechanisms, where agents iteratively record natural-language critiques of their previous actions, guiding future behavior to avoid recurring mistakes.
Additionally, AdaPlanner\citep{sun2023adaplanner} introduces closed-loop adaptive planning, allowing agents to refine their strategies on-the-fly based on environmental feedback, effectively reshaping action sequences in response to immediate outcomes. Similarly, Self-Refine\citep{madaan2023selfrefine} employs an iterative refinement loop in which the agent repeatedly critiques and revises its initial outputs, significantly improving task accuracy without explicit retraining. SICA (Self-Improving Coding Agent)\citep{robeyns2025selfimprovingcodingagent} further pushes the boundary by enabling agents to autonomously edit their underlying code and tools, iteratively enhancing their core reasoning abilities through direct self-modification.
From a reinforcement learning perspective, frameworks such as RAGEN\citep{wang2025ragen} and DYSTIL\citep{wang2025dystildynamicstrategyinduction} conceptualize multi-step tool-use tasks as Markov Decision Processes, optimizing agent policies through rich environmental rewards and strategy induction loops. RAGEN leverages dense feedback from the environment to iteratively fine-tune action policies, while DYSTIL utilizes high-level strategy advice generated by language models to progressively internalize complex decision-making skills into reinforcement learning agents.
Collectively, these approaches highlight a compelling paradigm where self-evolving agents not only leverage self-generated data but actively reshape their environments and internal mechanisms to fuel ongoing learning. Such dynamic interaction loops point toward autonomous, open-ended improvement cycles deeply grounded in experiential adaptation.


\subsection{Context}

An essential component of an LLM agent to be evolved is the context, which shapes how an agent behaves. To start with, we want to interpret two terms, "prompt optimization" and "memory evolution", which have been used in different literature. In most cases, these two terms can be used interchangeably because they both refer to what is included in the context window. Prompt optimization asks "how can we phrase or structure the instructions so the LLM behaves better?", and attends to details such as the wording, ordering. On the other hand, memory evolution asks "how should we store, forget, and retrieve context so that the agent can stay informed and perform better?", which focuses on what past information to surface or archive.

\subsubsection{Memory Evolution} 

LLM-based agents are increasingly designed with long-term memory mechanisms that grow and adapt as the agent continues to solve tasks and interacts with its environment\citep{shan2025cognitive, qian2023experiential}. An evolving memory enables the agent to accumulate knowledge, recall past events, and adjust its behavior based on experience. Many works stress that effective memory management is crucial for agent performance\citep{zhong2024memorybank, zhang2025g, yan2024efficient}. SAGE\citep{liang2024sage} uses the Ebbinghaus forgetting curve to decide what to remember or forget. A-mem\citep{xu2025mem} updates the agent memory structure to create interconnected knowledge networks through dynamic indexing and linking, following the basic principles of the Zettelkasten method. Mem0\citep{chhikara2025mem0} introduces a two-phase pipeline where the agent first extracts salient facts from recent dialogue and then decides how to update the long-term memory: the agent can ADD new facts, MERGE/UPDATE redundant ones, or DELETE contradictions. Furthermore, Memory-R1~\citep{yan2025memoryr1} presents a reinforcement learning framework to train a dedicated Memory Manager agent that learns to select structured operations like ADD, UPDATE, and DELETE.
Such a mechanism ensures the agent's long-term memory is coherent and up-to-date. MemInsight\citep{salama2025meminsight} augments raw memories with semantic structure, which summarizes and tags past interactions for retrieval later. REMEMBER\citep{zhang2024large} combines an LLM with a memory of experiences and uses reinforcement learning signals to decide how to update that memory after each episode. 
Memento~\citep{zhou2025memento} enables continual adaptation without fine-tuning the LLM's parameters by employing online reinforcement learning to optimize a case-retrieval policy, which allows the agent to learn from past experiences stored in an evolving memory bank. MemGen~\citep{zhang2025memgen} introduces a dynamic generative memory that operates in a latent space. It uses a learned memory trigger to decide when to invoke memory and a weaver to construct latent token sequences, enabling a fluid interweaving of reasoning and memory.


A critical aspect of memory evolution is enabling agents to learn heuristics or skills from past experiences. Rather than only retrieving exact past instances, advanced agents distill experiences into more general guidance\citep{zhao2024expel, fu2024autoguide}. Expel\citep{zhao2024expel} processes past trajectories to generate insights and rules to guide further interactions. This experiential knowledge accumulation leads to measurable gains, as the agent steadily performs better with more experience. ReasoningBank~\citep{ouyang2025reasoningbank} further develops this idea by distilling generalizable reasoning strategies from both successful and failed experiences into a structured memory. It also introduces memory-aware test-time scaling to generate diverse experiences on each task.
Other systems focus on storing higher-level building blocks of problem-solving. For instance, Agent Workflow Memory\citep{wang2024agent} records common sub-task sequences (workflows) so that an agent solving a complex task can retrieve and reuse a proven sequence of actions rather than plan from scratch. Similarly, MUSE~\citep{yang2025learningonthejob} introduces an experience-driven agent for long-horizon tasks, centered on a hierarchical memory module that organizes experience into strategic, procedural, and tool-use memories. The agent populates this memory through a Plan-Execute-Reflect-Memorize loop, enabling it to learn on the job.
In the Richelieu diplomacy agent, the system improves its negotiation strategies by augmenting its memory through self-play games, storing the insights from simulated interactions to refine future decisions\citep{guan2024richelieu}.
By generalizing from specific episodes to reusable knowledge, these approaches illustrate how memory evolution turns an agent’s one-time experiences into long-term competencies, which leads to agents evolving.

\subsubsection{Prompt Optimization}
While memory evolution focused on what knowledge an agent retains, Prompt Optimization (PO) enables LLM agents to self-evolve by refining the instructions it feeds to the backbone model, which directly alters the model's behavior without modifying model weights\citep{ramnath2025systematic}. Early research treats instruction design as a search problem. APE\citep{zhou2022large} generates candidate prompts, scores them on validation examples, and selects the best. ORPO\citep{yang2023large} extends this idea by letting the model iteratively rewrite its own prompt, guided by feedback on prior outputs. ADO\citep{lin2025ado} introduces DSP that imposes semantic constraints on iteratively proposed prompts to facilitate finding the optimal prompt. ProTeGi\citep{pryzant2023automatic} generates natural language "corrections" that are applied as edits to the prompt, forming a textual analogue of gradient descent. PromptAgent\citep{wang2023promptagent} casts prompt discovery as Monte-Carlo Tree Search, exploring instruction space strategically, while evolutionary approaches like PromptBreeder\citep{fernando2023promptbreeder} maintain a population to discover increasingly effective instructions. REVOLVE\citep{zhang2024revolve} further stabilizes long optimization runs by tracking the trajectory of model responses and applying smoothed updates. Pushing this autonomy to its limit, SPO\citep{xiang2025self} creates a fully self-contained loop where the model generates its training data and uses pairwise preference comparison on its outputs to refine the prompt, eliminating the need for any external labeled data or human feedback. Collectively, these techniques demonstrate that an agent can autonomously improve its prompting policy, turning prompt text into a learnable component that co-evolves with the agent’s experience. To address the brevity bias and context collapse of some optimizers, Agentic Context Engineering (ACE)~\citep{zhang2025ACE} treats contexts as comprehensive playbooks that accumulate strategies over time. It uses a modular agentic process with incremental updates to evolve contexts for both offline prompt optimization and online memory adaptation.


In complex systems, an agent often orchestrates a sequence of LLM calls or collaborates with other agents, making prompt design a multi-node problem. Frameworks such as DSPy represent an entire workflow as a graph whose sub-prompts are jointly tuned for a global objective\citep{khattab2023dspy}. Trace\citep{wang2023trace}, TextGrad\citep{yellamraju2024textgrad}, and LLM-AutoDiff\citep{yin2025llm} generalize this idea by treating each prompt as a parameter in a differentiable program and propagating natural-language “gradients” to refine every step. In collaborative scenarios, Multi-Agent System Search (MASS)\citep{zhou2025maasprompt} first optimizes individual role prompts and then refines inter-agent communication patterns, while MAS-ZERO\citep{ke2025maszero} dynamically proposes and revises role prompts to assemble an effective team for each new problem. Evolutionary systems such as EvoAgent\citep{yuan2024evoagent} and AgentSquare\citep{shang2025agentsquare} treat each agent along with prompts as the modules and use mutation and selection to discover specialized teams that outperform hand-crafted designs. These approaches extend PO from a single instruction to the language that defines whole workflows or societies of agents.



\subsection{Tools}


An agent's capabilities are fundamentally defined by the tools it can wield. The trajectory of agent development is marked by a crucial evolution: from being mere tool users to becoming autonomous tool makers. This transition from relying on predefined, static toolsets to enabling agents to autonomously expand and refine their own skills is a critical leap towards cognitive self-sufficiency. This paradigm, where agents dynamically adapt their capabilities, allows them to solve a long tail of complex problems not envisioned by their initial designers. This evolution unfolds across three interconnected fronts: tool discovery, mastery, and management, as detailed in the subsections below.

\paragraph{Autonomous Discovery and Creation}
The primary impetus for autonomous tool creation is to overcome the inherent limitations of a fixed toolset, granting agents the flexibility to innovate on demand. Methodologies for this now span a spectrum from opportunistic discovery to formalized synthesis. At one end, agents like Voyager build an ever-expanding library of skills through emergent trial-and-error, driven by an intrinsic motivation to explore complex, open-ended environments like Minecraft\citep{wang2023voyager}. This exploratory approach is powerful for generating a wide array of skills but may lack precision. In contrast, systems like ATLASS, Alita, and Live-SWE-Agent take a more reactive approach, often creating new tools from scratch or employing retrieval-augmented generation (RAG) to search open-source code repositories the moment a capability gap is identified\citep{Haque2025ATLASS, qiu2025alita, qiu2025alitag, xia2025livesweagent}.
At the other end of the spectrum lie highly structured frameworks that treat tool creation as a deliberate engineering process. CREATOR, for example, disentangles abstract tool creation (e.g., reasoning about the general structure of a reusable function for averaging temperatures over N days) from concrete tool usage (e.g., deciding how to apply that function to a specific city and time range), which enhances modularity and reusability\citep{qian2023creator}. Even more formally, SkillWeaver analyzes successful human or agent task trajectories to propose, synthesize, and hone new skills into robust, reusable APIs, ensuring a higher degree of initial quality\citep{zheng2025skillweaver}. Furthermore, frameworks like CRAFT demonstrate that creating specialized toolsets for specific domains is essential to complement general-purpose models, enabling expert-level performance without sacrificing adaptability\citep{yuan2024craft}. RL-GPT\citep{liu2024rl} integrates generated code implementations into the RL pipeline, leveraging these as tools to tackle complex tasks while addressing simpler ones directly using a Code-as-Policy approach. This integration dynamically adapts and evolves in response to environmental feedback, enabling continuous improvement. However, this burgeoning autonomy introduces significant challenges, particularly around safety and security. The unconstrained generation of code risks creating tools with exploitable vulnerabilities or unintended harmful behaviors, making automated verification and sandboxing critical areas for future research.

\paragraph{Mastery Through Iterative Refinement}
The proliferation of self-created tools necessitates a robust mechanism for their mastery; a newly generated tool is often a brittle script, not a reliable function. This is where iterative refinement becomes essential. Frameworks like LearnAct and From Exploration to Mastery establish a critical self-correction loop where the agent learns from its own experience\citep{zhao2024learnact, fromexplorationtomastery}. This involves tackling the difficult "credit assignment" problem: determining precisely which line of code or which parameter was responsible for a failure. To do this, the agent analyzes a rich variety of feedback signals—including compiler errors, unexpected API return values, environmental state changes, or even implicit signals from a user's subsequent actions. The goal is not only to debug the tool's underlying code but also to refine its documentation (e.g., its docstring and argument descriptions), which is crucial for improving the agent's ability to understand and correctly use the tool in the future.
This refinement process also opens the door for valuable human-agent collaboration. While full autonomy is the ultimate goal, many systems can be designed with a "human in the loop," where a human expert can provide corrections, offer high-level suggestions, or validate a newly created tool. This collaborative approach can significantly accelerate the mastery process and ensure that the agent's skills align with human intentions and safety standards. Ultimately, this self-honing process is what elevates a nascent skill into a dependable capability, ensuring the agent's growing skill library increases not just in quantity, but more importantly, in quality and robustness.

\paragraph{Scalable Management and Selection}
As an agent's mastered skill library grows into the hundreds or thousands, it faces a "curse of abundance." The challenge shifts from creating tools to efficiently managing and selecting from them. A large library creates a massive search space, making traditional retrieval methods slow and inaccurate. To overcome this, ToolGen~\citep{wang2025toolgen} represents a fundamental paradigm shift by encoding tools as unique tokens within the language model's vocabulary. This elegantly reframes tool retrieval as a generation problem, leveraging the transformer's immense pattern-recognition capabilities to predict the most appropriate tool as a natural continuation of its thought process. TOOLMEM~\citep{xiao2025toolmem} enables agents to learn and store the strengths and weaknesses of different tools in a dedicated memory. At inference, the agent retrieves this knowledge to make more informed decisions, optimizing tool selection for specific task requirements.
Beyond selecting a single tool, advanced agents must also excel at tool composition—learning to chain multiple tools in novel sequences to solve multi-step problems. This is a higher-order management task. Architectural approaches like AgentSquare engage in a form of meta-learning, automatically searching the modular design space of an agent—including its planning, memory, and tool-use components—to find an optimal configuration for complex task execution\citep{shang2025agentsquare}. As a logical endpoint to this evolutionary trend, visionary concepts like the Darwin Godel Machine propose a framework for open-ended evolution, where the agent can fundamentally rewrite its own core code. In this vision, the distinction between the agent and its tools blurs, leading to a recursive cascade of self-improvement that transcends tool enhancement alone\citep{zhang2025darwin}. In essence, this entire evolutionary path aims to establish a closed and virtuous cycle: a truly autonomous agent that can perceive gaps in its capabilities, create novel solutions, master them through practice, and seamlessly integrate them into a coherently managed and ever-expanding repertoire.




\subsection{Architecture}

The defining feature of next-generation agentic systems is their intrinsic capacity for self-improvement. This marks a fundamental shift from systems with fixed capabilities to those that can autonomously enhance their performance\citep{liu2025advances}. By treating their own internal logic and collaborative structures as optimizable components, these systems can adapt their behavior and design in response to feedback, achieving a level of efficiency and effectiveness that static designs cannot match. 
This section details how this self-optimization is realized, first by examining improvements within single-agent systems and then by exploring the co-evolution of complex multi-agent systems.



\subsubsection{Single-Agent System Optimization}

\paragraph{LLM-Invoking Node Optimization}

Optimizing a single LLM call is straightforward in isolation, but within an agentic system, it becomes a difficult credit assignment problem, as the effect of any single change is obscured by subsequent steps. Research addresses this by making node-level components optimizable, following two main strategies. The first focuses on refining nodes within a fixed agentic topology. A prime example is TextGrad \citep{yellamraju2024textgrad}, which, inspired by backpropagation, uses "textual gradients" to propagate feedback from the final output backward through the workflow, guiding systematic, local refinements at each node without altering the system's overall structure. The second, parallel strategy integrates this component-level optimization directly into the search for the system's architecture itself. Under this approach, node characteristics become tunable parameters in a larger search space. For instance, frameworks can embed prompt engineering directly into the search loop, allowing the system to discover not just the optimal workflow but also the most effective instruction for each agent simultaneously \citep{zhou2025maasprompt}. Similarly, EvoFlow \citep{zhang2025evoflow} uses evolutionary algorithms to construct heterogeneous workflows by selecting the most suitable LLM for each task from a diverse pool. This holistic strategy enables the discovery of systems that are co-optimized for both their structure and individual agent capabilities, effectively balancing metrics like overall performance and cost \citep{ye2025xmas}.

\paragraph{Autonomous-Agent Optimization}
Building upon the optimization of individual LLM-invoking nodes, a more profound level of self-improvement targets the autonomous agent as a holistic entity. This evolution proceeds along two main fronts: optimizing the agent's high-level architectural design and enabling the agent to directly modify its own source code. The first approach focuses on discovering the optimal agent structure. AgentSquare \citep{shang2025agentsquare} exemplifies this by defining a modular design space of components like planners and memory modules, then using an evolutionary algorithm to find the most effective combination for a given task. The second front involves agents that dynamically rewrite their own operational code. This is seen in radical systems like the Darwin Gödel Machine \citep{zhang2025darwin}, which recursively modifies its own Python codebase, and AlphaEvolve \citep{novikov2025alphaevolve}, which uses evolutionary coding to improve specific algorithms. Similarly, Gödel Agent \citep{yin2025godelagentselfreferentialagent} provides a self-referential framework for agents to analyze and alter their logic. Together, these two directions (optimizing the agent's architectural “blueprint” and its functional code) demonstrate a key trend toward turning the agent's fundamental structure and logic into learnable components. MemEvolve~\citep{zhang2025memevolve} introduces a meta-evolutionary framework that evolves not just the agent's experiential memory, but the memory system's architecture itself. Through a bilevel optimization process, it adapts the mechanisms for encoding, storing, and retrieving information to better suit specific task domains.

\subsubsection{Multi-Agent System Optimization}

How agents are organized and communicate within a system (its topology) fundamentally determines its capacity for solving complex problems. The field has evolved from using fixed, human-designed communication structures to creating dynamic systems that automatically adapt their organization to a given task, allowing them to discover and exploit the most effective collaboration patterns. This evolution is explored along two major fronts: the optimization of static, explicit workflows and the co-evolution of dynamic, internal policies.
\paragraph{Agentic Workflow Optimization}

The optimization of agentic workflows focuses on finding the most effective, often static, structure of communication and task delegation for a given problem. Early research established important foundations, with studies like AutoFlow \citep{li2024autoflow} demonstrating the automated creation of linear workflows from natural language, and GPTSwarm \citep{zhuge2024gptswarm} proposing a unifying graph-based framework. Concurrently, other foundational work explored how agents could evolve by using symbolic learning to distill their interaction experiences into an explicit, interpretable set of logical rules to guide future decisions \citep{zhou2024symbolic}. This abstraction of systems into tunable components—whether nodes, edges, or symbolic rules—was crucial. However, these early systems often lacked a formal method for efficiently navigating the vast space of possible configurations and interactions.

The major breakthrough came when ADAS \citep{hu2024automated} and AFlow \citep{zhang2024aflow} formally defined this challenge as a search and optimization problem. ADAS set a theoretical vision by framing system design as a search through a Turing-complete space of code-based configurations. Building on this, AFlow made it practical by introducing reusable operators that represent common agentic patterns and by employing Monte Carlo Tree Search (MCTS) to efficiently navigate the enormous design space.
Together, these works established a core methodology for treating agent system design as a tractable optimization problem, proving that automatically discovered workflows could outperform human-designed ones.

Following this formalization, research rapidly diversified toward creating customized agent systems for each specific query. Two primary strategies emerged: search-based and learning-based generation. Search-based methods, such as MaAS \citep{zhang2025maas}, create a "supernet" of potential architectures and then sample a specialized system from it. In parallel, learning-based methods train models to generate effective topologies directly. ScoreFlow \citep{wang2025scoreflow}, for instance, trains a generator using a novel preference optimization method, while FlowReasoner \citep{gao2025flowreasoner} uses reinforcement learning to train a meta-agent that constructs a bespoke workflow on the fly. This line of query-specific generation continues to be an active area of research \citep{ye2025masgpt, ke2025maszero}. Furthermore, it is important to note that this process is not limited to the topology alone; many of these frameworks also perform node-level optimization in tandem, such as co-optimizing prompts or selecting heterogeneous models as an integral part of the architectural generation process \citep{zhang2024aflow, zhou2025maasprompt, zhang2025evoflow}.

A key challenge for all search and learning methods is the computational cost of evaluating each potential workflow \citep{shang2025agentsquare}. To address this, researchers have developed lightweight prediction models. Agentic Predictor \citep{trirat2025agentic} is a prime example, training a model to accurately estimate a workflow's performance based on its structural and semantic features without a full execution. By providing a fast and inexpensive evaluation proxy, these predictors significantly accelerate the optimization process, making the exploration of vast design spaces feasible \citep{zhang2025gnns}.


\paragraph{Multi-Autonomous-Agent Optimization}
Distinct from optimizing a system's explicit workflow structure, this line of research focuses on how multiple autonomous agents can co-evolve their internal behavioral policies through interaction. This approach enables emergent capabilities like coordination, task delegation, and beneficial competition. For instance, ReMA \citep{wan2025remalearningmetathinkllms} uses multi-agent reinforcement learning (MARL) to collaboratively train a high-level meta-thinker and a low-level executor, significantly improving performance on reasoning benchmarks. Building on this, GiGPO \citep{feng2025groupingrouppolicyoptimizationllm} enhances MARL training by aggregating trajectories to provide more precise credit assignment, boosting success rates on long-horizon tasks. To support this direction, platforms like MARTI \citep{liao2025marftmultiagentreinforcementfinetuning} provide open-source infrastructure for orchestrating and scaling the training of these language-model collectives. Collectively, these studies underscore multi-agent reinforcement learning as a promising route for cultivating group-level competencies unattainable by individual agents alone.


